\documentclass[11pt]{article}
\usepackage{amsmath, amssymb, amsthm}
\usepackage{geometry}
\usepackage{booktabs}
\usepackage{xcolor}

\geometry{margin=1in}

\newcommand{\vect}[1]{\mathbf{#1}}
\newcommand{\mat}[1]{\mathbf{#1}}
\DeclareMathOperator{\diag}{diag}
\DeclareMathOperator{\tr}{tr}

\title{Exact Bias Correction for Intensity-Weighted Moments\\of Pedestal-Removed Gaussian Blobs}
\author{}
\date{}

\begin{document}
\maketitle

\section{Blob Rendering Model}

A dark blob on a white background is rendered as
\begin{equation}\label{eq:render}
  v(\vect{x}) = \frac{1 - G(\vect{x})}{1 - c}, \qquad
  G(\vect{x}) = \exp\!\Bigl(-\tfrac{1}{2}\,\vect{x}^\top \mat{\Sigma}^{-1}\vect{x}\Bigr),
\end{equation}
for $\|\vect{x}\|_{\mat{\Sigma}^{-1}} \le k$ (Mahalanobis radius), where
$c = e^{-k^2/2}$ is the Gaussian value at the measurement boundary and $k$ is
the \texttt{sigma\_cutoff} parameter. The denominator $1-c$ normalizes the
rendered intensity to $[0,1]$.

At the blob center $v = 0$ (dark); at the boundary $v = 1$ (white background).

\section{Inverted Intensity Profile}

The moment estimator uses the \emph{inverted} intensity
$f(\vect{x}) = B - v(\vect{x})$, where $B$ is the local background (maximum
pixel value in the measurement region, $B \approx 1$). This gives
\begin{equation}\label{eq:f}
  f(\vect{x}) = \frac{G(\vect{x}) - c}{1 - c},
\end{equation}
a pedestal-removed Gaussian that equals 1 at the center and vanishes at the
measurement boundary $r = k\sigma$.

\section{Measured Covariance}

The intensity-weighted second moment matrix is
\begin{equation}
  \mat{M} = \frac{\int_{D_k} \vect{x}\,\vect{x}^\top\, w(\vect{x})\,dA}
                  {\int_{D_k} w(\vect{x})\,dA},
\end{equation}
where $D_k = \{\vect{x} : \|\vect{x}\|_{\mat{\Sigma}^{-1}} \le k\}$ is the
measurement disc and $w(\vect{x})$ is the per-pixel weight function.

We consider two weight functions:
\begin{enumerate}
  \item \textbf{Non-spatial}: $w = f(\vect{x})$
  \item \textbf{Spatial Gaussian}: $w = f(\vect{x}) \cdot G(\vect{x})$
\end{enumerate}

In the isotropic case $\mat{\Sigma} = \sigma^2\mat{I}$, both reduce to 1D
integrals in polar coordinates. We derive
$\mat{M} = \alpha\,\sigma^2\,\mat{I}$ where $\alpha < 1$ is the bias factor,
and the correction is $\mat{\Sigma}_\text{true} = \alpha^{-1}\,\mat{M}$.

\subsection{Non-spatial case}

The weight $w = f = (G - c)/(1-c)$ decomposes into a Gaussian term and a
constant pedestal:
\begin{align}
  M_0 &= \int_{D_k} (G - c)\,dA = 2\pi\sigma^2(1-c) - c\pi k^2\sigma^2
       = \pi\sigma^2\bigl[2(1-c) - ck^2\bigr], \\[4pt]
  M_{xx} &= \int_{D_k} x^2(G - c)\,dA.
\end{align}

For the Gaussian term, using the substitution $u = r^2/(2\sigma^2)$:
\begin{equation}
  \int_{D_k} x^2\,G\,dA
    = \pi\!\int_0^{k\sigma}\!r^3 e^{-r^2/(2\sigma^2)}dr
    = 2\pi\sigma^4\bigl[1 - (1 + \tfrac{k^2}{2})\,c\bigr].
\end{equation}

For the constant pedestal over the disc:
\begin{equation}
  \int_{D_k} x^2\cdot c\,dA
    = c\pi\!\int_0^{k\sigma}\!r^3\,dr
    = \frac{c\pi k^4\sigma^4}{4}.
\end{equation}

The measured variance is therefore
\begin{equation}\label{eq:nonspatial}
  \boxed{
  \frac{\sigma^2_\text{meas}}{\sigma^2}
  = \frac{M_{xx}}{M_0\,\sigma^2}
  = \frac{2\bigl[1-(1+\frac{k^2}{2})c\bigr] - \frac{ck^4}{4}}
         {2(1-c) - ck^2}.
  }
\end{equation}

\subsection{Spatial Gaussian case}

With spatial weight $G(\vect{x})$, the total weight is
\begin{equation}
  w = f \cdot G = \frac{G^2 - cG}{1-c}.
\end{equation}

$G^2 = \exp(-r^2/\sigma^2)$ is an unnormalized Gaussian with effective
parameter $\sigma/\sqrt{2}$; $cG$ is $c$ times the original Gaussian with
parameter~$\sigma$. The truncated integrals over $D_k$ evaluate to:
\begin{align}
  M_0^{(G^2)} &= \pi\sigma^2\bigl(1 - e^{-k^2}\bigr), &
  M_{xx}^{(G^2)} &= \frac{\pi\sigma^4}{2}\bigl[1 - (1+k^2)e^{-k^2}\bigr], \\[4pt]
  M_0^{(cG)} &= 2c\pi\sigma^2(1 - c), &
  M_{xx}^{(cG)} &= 2c\pi\sigma^4\bigl[1 - (1+\tfrac{k^2}{2})c\bigr].
\end{align}

The measured variance is
\begin{equation}\label{eq:spatial}
  \boxed{
  \frac{\sigma^2_\text{meas}}{\sigma^2}
  = \frac{\frac{1}{2}[1-(1+k^2)e^{-k^2}] - 2c[1-(1+\frac{k^2}{2})c]}
         {[1-e^{-k^2}] - 2c(1-c)}.
  }
\end{equation}

\subsection{Correction factors}

The correction multiplier applied to the measured covariance is
$\gamma = \sigma^2/\sigma^2_\text{meas}$, the reciprocal of the bias factor.

\begin{table}[h]
\centering
\caption{Exact correction factors $\gamma$ vs.\ prior approximate formulas.}
\label{tab:corrections}
\begin{tabular}{@{}ccccc@{}}
\toprule
$k$ & Non-spatial (exact) & Non-spatial (old) & Spatial (exact) & Spatial (old) \\
\midrule
2.0 & 1.837 & 1.456 & 2.548 & 2.161 \\
2.5 & 1.355 & 1.168 & 2.171 & 2.024 \\
3.0 & 1.136 & 1.053 & 2.044 & 2.002 \\
3.5 & 1.044 & 1.014 & 2.009 & 2.000 \\
4.0 & 1.011 & 1.003 & 2.001 & 2.000 \\
\bottomrule
\end{tabular}
\end{table}

The old non-spatial formula was the standard Gaussian truncation correction
\begin{equation}
  \gamma_\text{old} = \frac{1 - e^{-k^2/2}}{1 - (1+\tfrac{k^2}{2})e^{-k^2/2}},
\end{equation}
which ignores the constant pedestal term. The old spatial formula was
$2 \times \gamma_\text{old}(k\sqrt{2})$, which treats the weight as a pure
$G^2$ without the $cG$ pedestal correction.

\section{Iterative Correction Placement}

\subsection{The catastrophic shrinkage bug}

The iterative moment estimator refines the measurement region from circular to
elliptical over multiple iterations. At each iteration, the estimated
$\mat{\Sigma}$ defines the Mahalanobis boundary for the next measurement.

If the bias correction $\gamma$ is applied only \emph{after} all iterations
(post-hoc), the raw biased $\mat{\Sigma}_\text{meas}$ is used to define
increasingly smaller measurement regions:
\begin{equation}
  \mat{\Sigma}_\text{meas}^{(t)} \approx \gamma^{-1}\,\mat{\Sigma}_\text{true}
  \implies \text{measurement region shrinks by } \gamma^{-1/2} \text{ per iteration.}
\end{equation}

For the spatial case at $k = 3$, $\gamma \approx 2.04$, so the raw measured
$\sigma$ is $\approx 0.70\,\sigma_\text{true}$. Each iteration compounds this
error, and after 5 iterations:
\begin{equation}
  \frac{\hat\sigma}{\sigma_\text{true}} \approx 0.47 \quad
  \text{(catastrophic collapse)}.
\end{equation}

For the non-spatial case ($\gamma \approx 1.14$ at $k=3$), the effect is
milder: the ratio degrades from $0.96$ (1 iteration) to $0.94$ (5 iterations).

\subsection{Fix: correct inside each iteration}

The correction must be applied to $\mat{\Sigma}_\text{meas}$ \emph{before}
constructing the next iteration's measurement ellipse:
\begin{equation}
  \mat{\Sigma}^{(t)}_\text{corrected} = \gamma\,\mat{\Sigma}^{(t)}_\text{meas},
  \qquad
  \mat{A}^{(t+1)} = k\,\bigl(\mat{\Sigma}^{(t)}_\text{corrected}\bigr)^{1/2}.
\end{equation}

This ensures the measurement region remains correctly sized across iterations.

\section{Experimental Validation}

\subsection{Isotropic synthetic blob}

A single isotropic blob with $\sigma_\text{true} = 10$ was rendered using the
pedestal-removed Gaussian model~\eqref{eq:render} and fitted with
\texttt{fit\_moments}.

\begin{table}[h]
\centering
\caption{Scale ratio $\hat\sigma/\sigma_\text{true}$ on isotropic synthetic blob.}
\label{tab:isotropic}
\begin{tabular}{@{}lccc@{}}
\toprule
Configuration & 1 iteration & 5 iterations (old) & 5 iterations (exact) \\
\midrule
Non-spatial & 0.964 & 0.945 & \textbf{1.001} \\
Spatial     & 0.996 & 0.467 & \textbf{1.009} \\
\bottomrule
\end{tabular}
\end{table}

\subsection{Anisotropic synthetic blob}

An anisotropic blob with $\sigma_a = 12$, $\sigma_b = 7$, $\theta = 30°$
(aspect ratio 1.714) was rendered and fitted.

\begin{table}[h]
\centering
\caption{Anisotropic blob recovery (5 iterations, exact correction).}
\label{tab:aniso}
\begin{tabular}{@{}lccc@{}}
\toprule
Method & $\hat\sigma_\text{eq}/\sigma_\text{eq}$ & Aspect ratio & $\hat\theta$ \\
\midrule
Non-spatial & 1.002 & 1.714 (exact) & $30.0°$ \\
Spatial     & 1.005 & 1.690 ($-1.4\%$) & $30.0°$ \\
Gradient tensor & 1.000$^*$ & 2.054 ($+20\%$) & $30.1°$ \\
\bottomrule
\multicolumn{4}{l}{\footnotesize $^*$Forced by $\det(\mat{\Sigma})=\sigma^4$ normalization.}
\end{tabular}
\end{table}

The gradient tensor overestimates anisotropy because the isotropic spatial
kernel $G_\sigma$ does not match the anisotropic blob shape, causing
directionally-dependent weighting of the gradient outer products.

\subsection{Rendered pattern (116 blobs)}

A $400\times400$ pattern was generated with blob scales $\sigma \in [3, 9.5]$.

\begin{table}[h]
\centering
\caption{Scale ratio statistics on rendered pattern (116 blobs, $k=3$).}
\label{tab:pattern}
\begin{tabular}{@{}lcc@{}}
\toprule
Method & Median $\hat\sigma/\sigma_\text{true}$ & Std \\
\midrule
Non-spatial & 0.988 & 0.001 \\
Spatial     & 0.996 & 0.004 \\
\bottomrule
\end{tabular}
\end{table}

Both methods recover aspect ratios very close to 1.0 (median 1.0006 and
1.0003 respectively) on the isotropic rendered blobs.

\section{Conclusions}

\begin{enumerate}
  \item The pedestal-removed Gaussian rendering model introduces a
        scale-dependent bias in intensity-weighted moments. The exact
        correction~\eqref{eq:nonspatial}--\eqref{eq:spatial} eliminates this
        bias for both spatial and non-spatial weighting.

  \item The truncation correction \emph{must} be applied inside each iteration
        of the Baumberg-style iterative refinement. Post-hoc correction causes
        catastrophic measurement region shrinkage, especially with spatial
        Gaussian weighting ($\gamma \approx 2$ at $k=3$).

  \item Spatial Gaussian weighting provides better contamination rejection from
        neighboring blobs and is the recommended mode for dense blob patterns.
        It achieves 0.4\% median scale error on rendered patterns.

  \item The gradient tensor (structure tensor) overestimates blob anisotropy by
        ${\sim}20\%$ when using an isotropic spatial kernel on anisotropic blobs.
        Intensity moments are the preferred shape estimator.
\end{enumerate}

\end{document}
